% from ega4-2-fr.tex
\oldpage[IV-2]{115}

Cela résulte de la forme (5.10.2, b)) de la condition pour qu'un
$\mathcal O_X$-Module quasi-cohérent soit $Z$-pur, et de (3.1.7).

\begin{corollary}[5.10.4]
\label{IV.2.s5.10.4-fr}
Supposons $\mathcal F$ cohérent. Pour que $\mathcal F$ soit $Z$-pur en un
point $x\in X$, il faut et il suffit que
\[
\operatorname{prof}_{Z_x}(\widetilde{\mathcal F}_x)\geq1
\]
(avec les notations de (5.9.6)).
\end{corollary}

Cela résulte aussitôt de (5.10.2) et (5.9.6).

\begin{theorem}[5.10.5]
\EGAReaderTarget{ega4.c04.thm.5-10-5}{IV.2.s5.10.5-fr}
Soient $X$ un préschéma localement noethérien, $Z$ une partie de $X$ stable
par spécialisation, $\mathcal F$ un $\mathcal O_X$-Module cohérent. Pour que
$\mathcal F$ soit $Z$-clos, il faut et il suffit que l'on ait
\[
\operatorname{prof}_Z(\mathcal F)\geq2.
\]
\end{theorem}

En vertu de (5.10.2), on peut se borner au cas où $\mathcal F$ est $Z$-pur
et $\operatorname{prof}_Z(\mathcal F)\geq1$. De plus, dire que
$\operatorname{prof}_Z(\mathcal F)\geq2$ équivaut à dire que pour toute partie
fermée $Z_\alpha$ de $Z$, $\operatorname{prof}_{Z_\alpha}(\mathcal F)\geq2$;
et de même, il résulte de (5.9.8) que dire que $\mathcal F$ est $Z$-clos
équivaut à dire que $\mathcal F$ est $Z_\alpha$-clos pour tout $\alpha$. On
peut donc déjà se borner au cas où $Z$ est fermé. La question étant locale,
il suffit, pour tout $x\in Z$, de prouver le théorème pour $\mathcal F|U$,
$U$ étant un voisinage ouvert affine de $x$, et l'on peut donc se borner au
cas où $X=U$ est affine. On sait alors que $\operatorname{Ass}(\mathcal F)$
est fini (3.1.6), et comme
$\operatorname{Ass}(\mathcal F)\subset X-Z$, il y a une section $f$ de
$\mathcal O_X$ au-dessus de $X$ telle que
$\operatorname{Ass}(\mathcal F)\subset X_f\subset X-Z$ (II, 4.5.4); on en
déduit que $f$ est $\mathcal F$-régulière (3.1.9) et que pour tout $y\in Z$,
on a $f_y\in\mathfrak m_y$, donc
\[
\operatorname{prof}(\mathcal F_y)
=1+\operatorname{prof}(\mathcal F_y/f_y\mathcal F_y)
\quad(0, 16.4.6).
\]
La condition $\operatorname{prof}_Z(\mathcal F)\geq2$ équivaut donc à
$\operatorname{prof}_Z(\mathcal F/f\mathcal F)\geq1$, ou encore (5.10.2) au
fait que $\mathcal F/f\mathcal F$ est $Z$-pur, et il suffit de voir que cette
dernière propriété équivaut au fait que $\mathcal F$ est $Z$-clos.

Considérons la suite exacte
\[
0\longrightarrow\mathcal F\xrightarrow{f}\mathcal F
\longrightarrow\mathcal F/f\mathcal F\longrightarrow0
\]
(l'homothétie de rapport $f:\mathcal F\to\mathcal F$ étant par hypothèse
injective); si on pose $W=X-Z$, on a le diagramme commutatif
\[
\xymatrix{
0\ar[r] & \Gamma(X,\mathcal F)\ar[r]^f\ar[d] &
\Gamma(X,\mathcal F)\ar[r]\ar[d] &
\Gamma(X,\mathcal F/f\mathcal F)\ar[r]\ar[d] & 0\\
0\ar[r] & \Gamma(W,\mathcal F)\ar[r]_f &
\Gamma(W,\mathcal F)\ar[r] & \Gamma(W,\mathcal F/f\mathcal F)
}
\]
où les deux lignes sont exactes ($X$ étant affine). Si l'homomorphisme de
restriction $\Gamma(X,\mathcal F)\to\Gamma(W,\mathcal F)$ est bijectif, on
déduit de ce diagramme que
\[
\Gamma(X,\mathcal F/f\mathcal F)
\longrightarrow\Gamma(W,\mathcal F/f\mathcal F)
\]
est injectif, et cela montre (5.9.8) que si $\mathcal F$ est $Z$-clos,
$\mathcal F/f\mathcal F$ est $Z$-pur. Inversement, supposons que
$\mathcal F/f\mathcal F$ soit $Z$-pur, et soit $s$ une section de $\mathcal F$
au-dessus de $W$; comme $X_f\subset W$, il existe un entier $n>0$ tel que
$f^n(s|X_f)$ se prolonge en une section $t$ de $\mathcal F$ au-dessus de $X$
(I, 1.4.1); d'ailleurs, les restrictions de $t$ et de $f^ns$ à $X_f$ étant
les mêmes, il en résulte que la restriction de $t$ à $W$ est égale à $f^ns$
en vertu de la relation $\operatorname{Ass}(\mathcal F)\subset X_f$
(5.10.2); comme $f$ est $\mathcal F$-régulier, il suffira de voir que $t$ est
de la forme $f^nt'$, où $t'\in\Gamma(X,\mathcal F)$

